\begin{textsample}{POS Dim 3 – ap – Score 15.00 – ap/ap\_inf\_10.txt}  \label{ex:f3_pos_142}
II – Campo de Experiência : Corpo, \textbf{Gestos} e Movimentos Com o corpo ( por meio dos sentidos, \textbf{gestos}, movimentos impulsivos ou intencionais, coordenados ou espontâneos ), as \textbf{crianças}, desde cedo, exploram o mundo, o espaço e os objetos do seu entorno, estabelecem relações, expressam -se, \textbf{brincam} e produzem conhecimentos sobre si, sobre o outro, sobre o universo social e cultural, tornando -se, \textbf{progressivamente}, conscientes dessa corporeidade.
Por meio das diferentes linguagens, como a música, a dança, o teatro, as \textbf{brincadeiras} de faz de conta, elas se comunicam e se expressam no entrelaçamento entre corpo, emoção e linguagem.
As \textbf{crianças} conhecem e reconhecem as sensações e funções de seu corpo e, com seus \textbf{gestos} e movimentos, identificam suas potencialidades e seus limites, desenvolvendo, ao mesmo tempo, a consciência sobre o que é seguro e o que pode ser um risco à sua integridade física.
Na Educação \textbf{Infantil}, o corpo das \textbf{crianças} ganha centralidade, pois ele é o partícipe privilegiado das práticas pedagógicas de \textbf{cuidado} físico, orientadas para a emancipação e a liberdade, e não para a submissão.
Assim, a instituição escolar precisa promover oportunidades ricas para que as \textbf{crianças} possam, sempre animadas pelo espírito \textbf{lúdico} e na interação com seus pares, explorar e vivenciar um amplo repertório de movimentos, \textbf{gestos}, olhares, \textbf{sons} e \textbf{mímicas} com o corpo, para descobrir variados modos de ocupação e uso do espaço com o corpo ( tais como sentar com apoio, rastejar, engatinhar, escorregar, caminhar \textbf{apoiando} -se em berços, mesas e cordas, saltar, escalar, equilibrar -se, correr, dar cambalhotas, alongar -se etc. ).
Direitos de aprendizagem com foco no campo de experiência corpo, \textbf{gesto} e movimentos : - CONVIVER com \textbf{crianças} e \textbf{adultos} e experimentar, de múltiplas formas, a gestualidade que marca sua cultura e está presente nos \textbf{cuidados} pessoais, dança, música, teatro, artes circenses, jogos, \textbf{escuta} de histórias e \textbf{brincadeiras}. - \textbf{BRINCAR}, utilizando movimentos para se expressar, explorar espaços, objetos e situações, \textbf{imitar}, jogar, imaginar, interagir e utilizar \textbf{criativamente} o repertório da cultura corporal e do movimento. - PARTICIPAR de diversas atividades de \textbf{cuidados} pessoais e do contexto social, de \textbf{brincadeiras}, encenações teatrais ou circenses, danças e músicas ; desenvolver práticas corporais e autonomia para cuidar de si, do outro e do ambiente. - EXPLORAR amplo repertório de movimentos, \textbf{gestos}, olhares, \textbf{sons} e \textbf{mímicas} ; descobrir modos de ocupação e de uso do espaço com o corpo e adquirir a compreensão do seu corpo no espaço, no tempo e no grupo. - EXPRESSAR corporalmente emoções, ideias e opiniões, tanto nas relações cotidianas como nas \textbf{brincadeiras}, dramatizações, danças, músicas, contação de histórias, dentre outras manifestações, empenhando -se em compreender o que outros também expressam. - CONHECER -SE nas diversas oportunidades de interações e explorações com seu corpo ; reconhecer e valorizar o seu pertencimento de gênero, étnico-racial e religioso.

% matched lemmas: adulto, apoiar, brincadeira, brincar, criança, criativamente, cuidado, escuta, gesto, imitar, infantil, lúdico, mímico, progressivamente, som
\end{textsample}
